\subsection{General Steps}
The plan is create the following sections:
\begin{enumerate}
    \item A general header file that loads packages that are fully or at least partially compatible with the \LaTeX\/ Tagging Project.
    \item Samples of text in single and multiple columns.  Some of these I will type myself.  For others I will test to see if the lipsum package can be used.
    \item Various equations and expressions with and with out array like environments.  These will include inline, displayed, and stacked equations.
    \item A variety of imported and generated images.  All of these will contain instances of alt text, actual text, or artifact tags.
    \item Test special packages like tcolorbox or mint to see if and how they work.
\end{enumerate}

After the creation of each section we will test the generated pdf using Blackboard Allie or Acrobat Accessibility Check.  And use one or more programs to convert the file to HTML. To facilitate this process, based on prior experience, this project has two shells that can be used for compilation, one with the document meta data and one without.

\subsection{Resources and Programs}

In creating this document we will reference a number of sources which include but are not limited to:
\begin{itemize}
    \item The \LaTeX\/ Tagging Project: \url{https://latex3.github.io/tagging-project/}
    \item AMS guidance for accessibility: \url{https://www.ams.org/accessibility/accessibility-guidance}
    \item \textit{Author Guidelines for Preparing Accessible Mathematics Content} produced through a collaboration of societies: \url{https://epubs.siam.org/pb-assets/author_guidelines_accessible_mathematics.pdf}
    \item Arxiv \textit{LaTeX Markup Best Practices for Successful HTML Papers}: \url{https://info.arxiv.org/help/submit_latex_best_practices.html}
\end{itemize}
By using these resources we should be able to make a pdf as accessible as it can be.  For examples of what this can sound like you can look here \url{https://latex3.github.io/tagging-project/documentation/wtpdf/fulldoc} at these complied examples from the \LaTeX\/ Tagging Project.

We will also explore using the following software in the conversion of files to HTML:
\begin{itemize}
    \item Pandoc Universal Document Converter: \url{https://pandoc.org/}
    \item \TeX 4ht: \url{https://tug.org/tex4ht/}
    \item \LaTeX ML: \url{https://math.nist.gov/~BMiller/LaTeXML/}
\end{itemize}
It is important to become familiar with these since HTML is the preferred format for making mathematics electronically accessible.